\section{The Emergence of Physics}
\label{sec:physics}

The model is experience-first, so physics is not the foundation but the large-scale limit of the note-dynamics.

\subsection{Experience is the foundation, not the top}

It is tempting to draw reality as a stack with a discrete computational substrate at the bottom and experience emerging at the top. That ordering is backward. Experience is the bottom, the only substance, what everything is made of. The discrete mathematics is a description of its form, layered on top. Read from the bottom, the first layer is experience itself, vibes noting vibes, all the way down, with no non-experiential floor beneath it. The next layer is computational form, a discrete rewriting network of the kind explored by Wolfram (2020), which is a mathematical abstraction of the structure and dynamics of the experiential substrate, not a meaningless machine underneath it. The bits are felt tones. The rewrites are note-dynamics. The computation is experience computing itself. Above that come conserved fields and currents at large scale, and finally observed physics, space and gravity and particles, as seen from far away.

\begin{proposition}[Form is not substance]
Water flowing in a pipe is described exactly by the Navier-Stokes equations, but water is not made of equations. The equations describe the form of the underlying molecular motion. Temperature is the same, real and useful yet a description of molecular motion, not a separate substance. The discrete mathematics stands to experience as Navier-Stokes stands to molecules, a description of the form of the substrate, layered on top, not the substrate itself.
\end{proposition}

So the correction to digital physics is this. The universe is not made of meaningless bits with experience added on top. It is made of experience, and the discrete computation is the abstract form read off it. Physics is recovered from experience. Experience is never recovered from physics.

\subsection{The candidate correspondences}

Each piece of physics is meant to emerge as a large-scale limit. Space and distance emerge as note-distance (Section~\ref{sec:geometry}). Spacetime and the light cone emerge as the note-metric together with the causal order (Section~\ref{sec:time}). Energy emerges as the flux of causal notes across a region, in the manner of Wolfram (2020). Curvature and gravity emerge as the combinatorial curvature of the notes. The Dirac equation and spin emerge as a quantum-walk coin, the graded relational-vibe complex of Section~\ref{sec:matter}, in the manner of the quantum cellular automata of Lent et al. (1993) and Watrous (1995). Conserved currents emerge as invariants of the rule. None of these is yet a derivation. Each is a correspondence whose limit must be proven, a long program. But the shapes line up, and discreteness places the model in the family of digital-physics programs that already exhibit these emergences, such as those of Fredkin (2003) and 't Hooft (2016).

\subsection{What is established, and what would distinguish it}

The honest summary is that the model is internally coherent and constructively specified, with a clear build target, but it is not yet an empirically tested theory, and its connection to established physics is a program rather than a result. The pieces that follow from the axiom and constraints are derived. The minimal-model choices, namely ternary tone, the eternal-expansion fate, hyperbolic geometry, structural attention, and asynchronous scheduling, are committed working choices with named alternatives held open. The harmony-as-energy and free-energy-as-mind readings are well-motivated candidates. The matter construction is open.

Because the model is fully discrete, it is fully simulable, so it can make coarse, checkable predictions. Agency above a coherence threshold predicts a sharp phase transition at a computable critical coherence, not a smooth ramp. The settling-time reading of deliberation predicts a measurable relation between a system's integration and its decision latency. The model differs from Integrated Information Theory (Tononi, 2008) in treating integrated information as a coherence diagnostic on an experiential substrate rather than the source of experience, so a case where the two disagree about which system is one mind would discriminate them. A genuinely discrete base predicts a smallest scale and small departures from perfect continuity, the standard signatures shared with causal-set programs (Bombelli et al., 1987).
